\documentclass[a4paper,12pt]{article}
\usepackage[T2A]{fontenc}
\usepackage[utf8]{inputenc}
\usepackage{cmap}
\usepackage[english, russian]{babel}
\usepackage{wrapfig}
\usepackage[backend=biber,style=gost-numeric]{biblatex} % стиль ГОСТ
\addbibresource{literatura.bib}
\usepackage{misccorr} % в заголовках появляется точка, но при ссылке на них ее нет
\usepackage{amssymb,amsfonts,amsmath,amsthm}  
\usepackage{indentfirst}
\usepackage[usenames,dvipsnames]{color} 
\usepackage[unicode,hidelinks]{hyperref}
% \hypersetup{%
%     pdfborder = {0 0 0}
\usepackage{multirow}
\usepackage{makecell,multirow} 
\usepackage{ulem}
%\usepackage{hyperref}
%\usepackage{graphicx,wrapfig}
%\renewcommand\epsilon{\varepsilon}
% \renewcommand\epsilon{\varepsilon}
%\graphicspath{{img/}}
\usepackage{geometry}
\geometry{left=1cm,right=1cm,top=2cm,bottom=2cm,bindingoffset=0cm,headheight=15pt}
\usepackage{fancyhdr} 
\linespread{1.2} 
\frenchspacing 
\renewcommand{\labelenumii}{\theenumii)} 
% \usepackage{caption}
%%%%%%%%%%%%%%%%%%%%%%%%%%%%%%%%%%%%%%%%%%%%%%%%%%%%%%%%%%%%%%%%%%%%%%%%%%%%%%%
%%%%%%%%%%%%%%%%%%%%%%%%%%%%%%%%%%%%%%%%%%%%%%%%%%%%%%%%%%%%%%%%%%%%%%%%%%%%%%%

\def\labauthor{Сарафанов И.Г.}
\def\labauthors{\labauthor}
\def\labnumber{127}
\def\labtheme{Таблица производных и интегралов}

%%%%%%%%%%%%%%%%%%%%%%%%%%%%%%%%%%%%%%%%%%%%%%%%%%%%%%%%%%%%%%%%%%%%%%%%%%%%%%%
	%применим колонтитул к стилю страницы
\pagestyle{fancy} 
	%очистим "шапку" страницы
\fancyhead{\LaTeX} 
	%слева сверху на четных и справа на нечетных
\fancyhead[L]{\labauthors} 
	%справа сверху на четных и слева на нечетных
\fancyhead[R]{Иродов \textnumero 1.186} 
	%очистим "подвал" страницы
\fancyfoot{} 
	% номер страницы в нижнем колинтуле в центре
\fancyfoot[C]{\thepage} 
%%%%%%%%%%%%%%%%%%%%%%%%%%%%%%%%%%%%%%%%%%%%%%%%%%%%%%%%%%%%%%%%%%%%%%%%%%%%%%%
\usepackage{float}
\usepackage[mode=buildnew]{standalone}
\usepackage{tikz} 
% \usepackage{subcaption}
\usepackage{tikz,csvsimple,physics}
\makeatother
%\usepackage{booktabs}
\usepackage{pgfplots, pgfplotstable}
\usepackage[outline]{contour}
\usepackage{tocloft}
\renewcommand{\cftsecleader}{\cftdotfill{\cftdotsep}}


\begin{document}
\begin{center}
\subsection*{Домашняя работа \textnumero 5}
\end{center}

\textbf{Условие:} Частица 1 ударяет покоящуюся частицу 2. После удара они разлетаются симметрично под углом $60^\circ$ между направлениями. Найти отношение масс.

\bigskip

\textbf{Решение}
Импульс системы сохраняется, так как внешние силы пренебрежимо малы во время короткого удара:

\[
\vec P_\text{до} = \vec P_\text{после}
\]

До удара:

\[
P_y = 0
\]

После удара одна частица отклоняется вверх, другая вниз, угол между направлениями $60^\circ$ $\Rightarrow$ каждая отклоняется на $30^\circ$ от оси $x$  

\vspace{2em}
Применим \textit{ЗСИ} по оси $y$:
\[
m_1 v_1 \sin 30^\circ - m_2 v_2 \sin 30^\circ = 0
\]

Отсюда:

\[
m_1 v_1 = m_2 v_2
\]

Применим \textit{ЗСИ} по оси $x$:

\[
m_1 v_0 = m_1 v_1 \cos 30^\circ + m_2 v_2 \cos 30^\circ
\]

Подставляем $m_2 v_2 = m_1 v_1$:

\[
m_1 v_0 = m_1 v_1 (\cos 30^\circ + \cos 30^\circ) = 2 m_1 v_1 \cos 30^\circ
\]

\[
v_1 = \frac{v_0}{2 \cos 30^\circ} = \frac{v_0}{\sqrt{3}}
\]

\[
v_2 = \frac{m_1}{m_2} v_1 = \frac{v_1 m_1}{m_2}
\]

Применим \textit{ЗСМЭ:}

\[
\frac{1}{2} m_1 v_0^2 = \frac{1}{2} m_1 v_1^2 + \frac{1}{2} m_2 v_2^2
\]

Подставляем $v_2 = \frac{m_1}{m_2} v_1$:

\[
\frac{1}{2} m_1 v_0^2 = \frac{1}{2} m_1 v_1^2 + \frac{1}{2} m_2 \left(\frac{m_1}{m_2} v_1\right)^2 = \frac{1}{2} m_1 v_1^2 + \frac{1}{2} \frac{m_1^2}{m_2} v_1^2
\]

\[
m_1 v_0^2 = v_1^2 \left( m_1 + \frac{m_1^2}{m_2} \right) = v_1^2 m_1 \left( 1 + \frac{m_1}{m_2} \right)
\]

\[
\frac{v_0^2}{v_1^2} = 1 + \frac{m_1}{m_2} \quad \Rightarrow \quad \frac{v_0^2}{(v_0^2/3)} = 3
\]

\[
1 + \frac{m_1}{m_2} = 3 \quad \Rightarrow \quad \frac{m_1}{m_2} = 2
\]

\[
\boxed{\frac{m_1}{m_2} = 2}
\]
\end{document}